% Gaussian vs t-Copula: Capturing Joint Crashes — research-note PDF.
%
% PROSE AND CODE ARE VERBATIM from the website article:
%   site/src/content/articles/gaussian-vs-t-copula.tsx
%   site/src/content/articles/t-copula-code.ts
% Metadata from lib/articles.ts; project-card fallbacks from lib/topics.ts.
%
% NO CHARTS, deliberately. The article's header records it as unexecuted: no
% data module, no computed numbers, prose and code only. Section 4.2 asserts
% the headline comparison without showing it, and that is the article's own
% choice — plotting it here would mean inventing results.
\documentclass{pyportfolios-note}

\noteshorttitle{Gaussian vs t-Copula}

\begin{document}
\thispagestyle{notecover}

\notemasthead{Quantitative Insights}{Quant Insights}{Q3 2026}{}

\notetitle{Gaussian vs t-Copula:\\Capturing Joint Crashes}

\notedeck{Two models of ``moving together.'' Only one believes in contagion.}

\textbf{The claim we test:} a correlation number tells you how two assets
co-move \textit{on average} — it says nothing about whether they crash
\textit{together}. The Gaussian copula assumes joint extremes are vanishingly
rare (zero tail dependence); the t-copula assumes crashes cluster. Using two
sector ETFs — Financials (XLF) and Technology (XLK) — through the GFC, COVID
and 2022, we measure which model matches how markets actually behave on the
worst days.

\notecard
  {Risk Management}
  {SciPy · NumPy · seaborn}
  {XLF vs XLK (SPDR sectors)}
  {Jan 2007 – Dec 2024}
  {Choosing a dependence model that doesn't underestimate crash correlation
   (2008's lesson)}

\begin{notecode}
import numpy as np
import pandas as pd
import matplotlib.pyplot as plt
import seaborn as sns
from scipy import stats

np.random.seed(42)
plt.rcParams["figure.dpi"] = 110
sns.set_style("whitegrid")
\end{notecode}

% ==================================================== 1 The disagreement =====
\noteneedroom{160bp}
\notesection{1.\notenumsep The disagreement}

A copula separates \textit{how assets co-move} from \textit{how each behaves
alone}. Two copulas can share the same correlation $\rho$ yet disagree
completely about the tails:

\begin{multicols}{2}
\begin{notebullets}
\item \textbf{Gaussian copula:} as moves get more extreme, the dependence
      \textit{fades}. Joint crashes are essentially coincidences. Lower-tail
      dependence $\lambda$ = \textbf{0}, always.
\item \textbf{t-copula:} dependence \textit{persists} into the extremes —
      crashes arrive together. $\lambda > 0$, controlled by the degrees of
      freedom $\nu$ (smaller $\nu$ = fatter joint tails).
\end{notebullets}

This isn't academic. The Gaussian copula priced the CDOs of 2008 — and its
zero-tail-dependence assumption is precisely what failed when housing defaults
arrived together.
\end{multicols}

% ====================================================== 2 The contenders =====
\noteneedroom{190bp}
\notesection{2.\notenumsep The contenders}

% The t-copula column carries a formula that cannot line-break, so it gets the
% width; 115 + 60 + 330 + 24bp of column separation = 529bp of the 532.8 measure.
\begin{notetable}{@{}p{115bp}p{60bp}p{330bp}@{}}
\thcell{} & \thcell{Gaussian copula} & \thcell{t-copula} \\
\theadrule
Parameters & $\rho$ & $\rho, \nu$ (degrees of freedom) \\ \rowrule
Lower-tail dependence $\lambda$ & 0 — always &
  $2\,t_{\nu+1}\!\big(-\sqrt{\tfrac{(\nu+1)(1-\rho)}{1+\rho}}\big) > 0$ \\ \rowrule
Joint crashes & Coincidences & Expected \\ \rowrule
As $\nu \to \infty$ & — & Converges to Gaussian \\
\end{notetable}

The t-copula \textit{contains} the Gaussian as a limit — so if the data prefers
small $\nu$, that's the data voting for tail dependence.

% ============================================================= 3 The test ====
\noteneedroom{150bp}
\notesection{3.\notenumsep The test}

\notesubsection{3.1\notenumsep Data: XLF vs XLK through three crises}

\begin{notecode}
TICKERS = ["XLF", "XLK"]
START, END = "2007-01-01", "2024-12-31"

def load_prices(tickers, start, end):
    """Adjusted-close prices: yfinance first, Stooq as fallback."""
    try:
        import yfinance as yf
        df = yf.download(tickers, start=start, end=end, auto_adjust=True, progress=False)["Close"]
        if not df.empty:
            return df[tickers].dropna()
    except Exception as exc:
        print(f"yfinance failed ({exc}); trying Stooq…")
    cols = {}
    for t in tickers:
        url = f"https://stooq.com/q/d/l/?s={t.lower()}.us&i=d"
        cols[t] = pd.read_csv(url, parse_dates=["Date"], index_col="Date")["Close"].rename(t)
    return pd.concat(cols, axis=1).loc[start:end].dropna()

px = load_prices(TICKERS, START, END)
rets = px.pct_change().dropna()
n = len(rets)
print(f"{n} trading days, {px.index[0].date()} → {px.index[-1].date()}")
print(f"Linear correlation: {rets['XLF'].corr(rets['XLK']):.2f}")
\end{notecode}

\notesubsection{3.2\notenumsep Strip the marginals: pseudo-observations}

Copulas work on ranks. Converting each series to its empirical percentile
removes the marginal distributions and leaves only the dependence structure.

\begin{notecode}
u = rets.rank() / (n + 1)          # pseudo-observations in (0,1)

fig, axes = plt.subplots(1, 2, figsize=(11, 4.6))
axes[0].scatter(rets["XLF"], rets["XLK"], s=4, alpha=0.4, color="steelblue")
axes[0].set_title("Raw daily returns"); axes[0].set_xlabel("XLF"); axes[0].set_ylabel("XLK")
axes[1].scatter(u["XLF"], u["XLK"], s=4, alpha=0.4, color="darkorange")
axes[1].set_title("Pseudo-observations (the copula's view)")
axes[1].set_xlabel("XLF percentile"); axes[1].set_ylabel("XLK percentile")
plt.tight_layout(); plt.show()

print("Note the crowding in the corners of the right plot — especially bottom-left (joint crashes).")
\end{notecode}

\notesubsection{3.3\notenumsep Fit both copulas}

Gaussian: correlation of normal scores. t-copula: same $\rho$ (via Kendall's
$\tau$), with $\nu$ chosen by maximum likelihood.

\begin{notecode}
# --- Gaussian copula fit ---
z = stats.norm.ppf(u.values)
rho_g = np.corrcoef(z.T)[0, 1]

# --- t-copula fit: rho from Kendall's tau, nu by profile likelihood ---
tau = stats.kendalltau(u["XLF"], u["XLK"]).statistic
rho_t = np.sin(np.pi * tau / 2)

def t_copula_loglik(nu, u, rho):
    x = stats.t.ppf(u, df=nu)
    cov = np.array([[1, rho], [rho, 1]])
    num = stats.multivariate_t(loc=[0, 0], shape=cov, df=nu).logpdf(x)
    den = stats.t.logpdf(x, df=nu).sum(axis=1)
    return np.sum(num - den)

nus = np.arange(2, 31)
lls = [t_copula_loglik(v, u.values, rho_t) for v in nus]
nu_hat = nus[int(np.argmax(lls))]

fig, ax = plt.subplots(figsize=(9, 3.8))
ax.plot(nus, lls, "o-", color="steelblue")
ax.axvline(nu_hat, color="crimson", ls="--", label=f"MLE: ν = {nu_hat}")
ax.set_xlabel("Degrees of freedom ν"); ax.set_ylabel("Copula log-likelihood")
ax.set_title("The data votes: profile likelihood over ν"); ax.legend()
plt.tight_layout(); plt.show()

print(f"Gaussian copula: rho = {rho_g:.3f}")
print(f"t-copula:        rho = {rho_t:.3f}, nu = {nu_hat}")
print("\nSmall ν = strong tail dependence. ν above ~30 would mean 'basically Gaussian'.")
\end{notecode}

% ======================================================= 4 The scoreboard ====
\noteneedroom{150bp}
\notesection{4.\notenumsep The scoreboard}

\notesubsection{4.1\notenumsep Tail dependence: the headline number}

Probability that one ETF is in its worst q\% \textit{given} the other is — as
q $\rightarrow$ 0.

\begin{notecode}
# closed-form lower-tail dependence of the t-copula
lam_t = 2 * stats.t.cdf(-np.sqrt((nu_hat + 1) * (1 - rho_t) / (1 + rho_t)), df=nu_hat + 1)
print(f"Lower-tail dependence λ — Gaussian: 0.000 (by construction)")
print(f"Lower-tail dependence λ — t-copula: {lam_t:.3f}")
\end{notecode}

\notesubsection{4.2\notenumsep Empirical test: joint crash frequency}

Count the days both ETFs landed in their worst 5\% (and 1\%) simultaneously —
then ask each fitted copula how many such days it \textit{predicts} over the
same sample size.

\begin{notecode}
def simulate_gaussian(rho, size):
    cov = np.array([[1, rho], [rho, 1]])
    z = np.random.multivariate_normal([0, 0], cov, size)
    return stats.norm.cdf(z)

def simulate_t(rho, nu, size):
    cov = np.array([[1, rho], [rho, 1]])
    g = np.random.multivariate_normal([0, 0], cov, size)
    chi = np.random.chisquare(nu, size) / nu
    x = g / np.sqrt(chi)[:, None]
    return stats.t.cdf(x, df=nu)

N_SIM = 2_000_000
u_g = simulate_gaussian(rho_g, N_SIM)
u_t = simulate_t(rho_t, nu_hat, N_SIM)

print(f"{'q':>4} {'empirical':>10} {'Gaussian':>10} {'t-copula':>10}   (joint-crash days per sample of {n})")
for q in [0.05, 0.01]:
    emp = float(((u['XLF'] < q) & (u['XLK'] < q)).mean())
    pg  = float(((u_g[:,0] < q) & (u_g[:,1] < q)).mean())
    pt  = float(((u_t[:,0] < q) & (u_t[:,1] < q)).mean())
    print(f"{q:>4.0%} {emp*n:>10.1f} {pg*n:>10.1f} {pt*n:>10.1f}")
\end{notecode}

\begin{notecode}
# visual: simulated pseudo-observations vs reality, tails highlighted
fig, axes = plt.subplots(1, 3, figsize=(13.5, 4.4))
sets = [(u.values, "Real data (XLF/XLK)", "black"),
        (u_g[:len(u)], f"Gaussian (ρ={rho_g:.2f})", "steelblue"),
        (u_t[:len(u)], f"t-copula (ρ={rho_t:.2f}, ν={nu_hat})", "crimson")]
for ax, (dat, title, c) in zip(axes, sets):
    ax.scatter(dat[:, 0], dat[:, 1], s=3, alpha=0.35, color=c)
    ax.axvline(0.05, color="gray", ls=":", lw=0.8); ax.axhline(0.05, color="gray", ls=":", lw=0.8)
    joint = ((dat[:, 0] < 0.05) & (dat[:, 1] < 0.05)).sum()
    ax.set_title(f"{title}\njoint 5% tail: {joint} days")
    ax.set_xlabel("XLF percentile")
axes[0].set_ylabel("XLK percentile")
plt.tight_layout(); plt.show()
\end{notecode}

% ============================================================= 5 Verdict =====
\noteneedroom{280bp}
\notesection{5.\notenumsep Verdict}

\textbf{The t-copula wins where it counts — the corner where portfolios die.}
Typical result on this pair: the Gaussian copula \textit{underpredicts} joint
1\%-tail days by a large factor, while the t-copula lands close to the observed
count. The likelihood agrees: the fitted $\nu$ is small, which is the data
explicitly rejecting the Gaussian limit.

\begin{multicols}{2}
\textbf{Why it matters:}

\begin{notebullets}
\item \textbf{Risk models built on Gaussian dependence} understate exactly the
      events that cause maximum damage — simultaneous crashes across holdings.
      Your 99\% ``diversified'' VaR is too optimistic in the only scenario you
      bought diversification for.
\item \textbf{The 2008 lesson, quantified} — the Gaussian copula's
      $\lambda = 0$ is the mathematical form of ``national housing markets
      can't all fall together.'' They did. Any dependence model with zero tail
      dependence repeats that mistake in new clothing.
\end{notebullets}

\textbf{The honest caveats:}

\begin{notebullets}
\item The t-copula is \textit{symmetric} — it implies joint booms are as likely
      as joint crashes. Equity data usually shows \textit{asymmetric}
      dependence (crashes cluster more than rallies). A skewed-t or Clayton
      copula captures that; the t is the pragmatic middle ground.
\item One $\nu$ for all seasons — dependence itself shifts with regimes (see
      our 60/40 case study). Time-varying copulas exist but multiply
      complexity.
\end{notebullets}
\end{multicols}

\textbf{Where to go next:} the \textbf{Copulas \& Tail Dependence} tutorial for
foundations; \textbf{CVaR / Expected Shortfall} for the risk measure that
consumes these joint-tail probabilities; and the CDO story in our curriculum's
credit module for the trillion-dollar version of this comparison.

\end{document}
